\textbf{Phần 1. Trắc nghiệm lựa chọn}
\setcounterpageref{ex}{0}

\begin{ex} %[5H5V28]
Một bộ xếp hình Tangram gồm 7 miếng ghép được cắt ra từ một miếng bìa hình vuông có cạnh dài $12\text{ cm}$. Tổng diện tích của cả 7 miếng ghép trong bộ Tangram đó là
\choice
{$48\text{ cm}^2$}
{\True $144\text{ cm}^2$}
{$72\text{ cm}^2$}
{$96\text{ cm}^2$}
\loigiai{
Tổng diện tích 7 miếng ghép bằng chính diện tích miếng bìa hình vuông ban đầu:\\
$12 \times 12 = 144\text{ (cm}^2\text{)}$.\\
\textbf{Đáp án của câu hỏi là B}
}
\end{ex}

\begin{ex} %[5H5V28]
Trong bộ Tangram được tạo từ hình vuông có diện tích $144\text{ cm}^2$, diện tích của miếng ghép hình vuông nhỏ bằng $\dfrac{1}{8}$ diện tích cả bộ. Diện tích của miếng ghép hình vuông nhỏ đó là
\choice
{$12\text{ cm}^2$}
{$16\text{ cm}^2$}
{\True $18\text{ cm}^2$}
{$24\text{ cm}^2$}
\loigiai{
Diện tích của miếng ghép hình vuông nhỏ là:
$144 \times \dfrac{1}{8} = 18\text{ (cm}^2\text{)}$.\\
\textbf{Đáp án của câu hỏi là C}
}
\end{ex}

\begin{ex} %[5H5V28]
Để cắt dán một bức tranh chiếc thuyền buồm, bạn Mai dùng một hình thang làm thân thuyền (hai đáy $15\text{ cm}$ và $25\text{ cm}$, chiều cao $8\text{ cm}$) và 2 hình tam giác làm cánh buồm (mỗi cánh buồm có đáy $10\text{ cm}$, chiều cao $14\text{ cm}$). Tổng diện tích giấy màu bạn Mai đã dùng là
\choice
{$230\text{ cm}^2$}
{\True $300\text{ cm}^2$}
{$320\text{ cm}^2$}
{$440\text{ cm}^2$}
\loigiai{
Diện tích thân thuyền là:
$(15 + 25) \times 8 : 2 = 160\text{ (cm}^2\text{)}$.\\
Diện tích hai cánh buồm là:
$2 \times (10 \times 14 : 2) = 140\text{ (cm}^2\text{)}$.\\
Tổng diện tích giấy màu cần dùng là:
$160 + 140 = 300\text{ (cm}^2\text{)}$.\\
\textbf{Đáp án của câu hỏi là B}
}
\end{ex}

\begin{ex} %[5H5V28]
Một chiếc diều hình thoi có độ dài hai đường chéo là $60\text{ cm}$ và $80\text{ cm}$. Diện tích của chiếc diều đó là
\choice
{$4\,800\text{ cm}^2$}
{\True $2\,400\text{ cm}^2$}
{$1\,200\text{ cm}^2$}
{$280\text{ cm}^2$}
\loigiai{
Diện tích hình thoi bằng nửa tích hai đường chéo:
$60 \times 80 : 2 = 2\,400\text{ (cm}^2\text{)}$.\\
\textbf{Đáp án của câu hỏi là B}
}
\end{ex}

\begin{ex} %[5H5V28]
Một bồn hoa có hình dạng gồm một hình vuông cạnh $4\text{ m}$ và bốn phía gắn thêm 4 nửa hình tròn có đường kính bằng cạnh hình vuông (đường kính $4\text{ m}$). Diện tích toàn bộ bồn hoa đó là
\choice
{$28{,}56\text{ m}^2$}
{\True $41{,}12\text{ m}^2$}
{$66{,}24\text{ m}^2$}
{$53{,}68\text{ m}^2$}
\loigiai{
Diện tích hình vuông là:
$4 \times 4 = 16\text{ (m}^2\text{)}$.\\
Bốn nửa hình tròn ghép lại thành 2 hình tròn đầy đủ có bán kính $r = 4 : 2 = 2\text{ (m)}$.\\
Diện tích 2 hình tròn là:
$2 \times (3{,}14 \times 2 \times 2) = 25{,}12\text{ (m}^2\text{)}$.\\
Tổng diện tích bồn hoa là:
$16 + 25{,}12 = 41{,}12\text{ (m}^2\text{)}$.\\
\textbf{Đáp án của câu hỏi là B}
}
\end{ex}

\begin{ex} %[5H5V28]
Để làm một bia bắn tên, người ta vẽ 2 đường tròn đồng tâm có bán kính lần lượt là $4\text{ cm}$ và $6\text{ cm}$. Diện tích của vòng tròn ngoài cùng (phần nằm giữa hai đường tròn) là
\choice
{$31{,}4\text{ cm}^2$}
{\True $62{,}8\text{ cm}^2$}
{$113{,}04\text{ cm}^2$}
{$50{,}24\text{ cm}^2$}
\loigiai{
Diện tích hình tròn lớn bán kính $6\text{ cm}$ là:
$3{,}14 \times 6 \times 6 = 113{,}04\text{ (cm}^2\text{)}$.\\
Diện tích hình tròn bé bán kính $4\text{ cm}$ là:
$3{,}14 \times 4 \times 4 = 50{,}24\text{ (cm}^2\text{)}$.\\
Diện tích phần vòng ngoài là:
$113{,}04 - 50{,}24 = 62{,}8\text{ (cm}^2\text{)}$.\\
\textbf{Đáp án của câu hỏi là B}
}
\end{ex}

\begin{ex} %[5H5V28]
Xung quanh một sân cỏ hình chữ nhật dài $20\text{ m}$, rộng $15\text{ m}$, người ta lát một lối đi rộng $1\text{ m}$ bao quanh phía ngoài sân cỏ. Diện tích của lối đi đó là
\choice
{$70\text{ m}^2$}
{\True $74\text{ m}^2$}
{$76\text{ m}^2$}
{$80\text{ m}^2$}
\loigiai{
Chiều dài toàn bộ khu đất (gồm cả lối đi) là:
$20 + 1 + 1 = 22\text{ (m)}$.\\
Chiều rộng toàn bộ khu đất (gồm cả lối đi) là:
$15 + 1 + 1 = 17\text{ (m)}$.\\
Diện tích toàn bộ khu đất là:
$22 \times 17 = 374\text{ (m}^2\text{)}$.\\
Diện tích sân cỏ là:
$20 \times 15 = 300\text{ (m}^2\text{)}$.\\
Diện tích lối đi là:
$374 - 300 = 74\text{ (m}^2\text{)}$.\\
\textbf{Đáp án của câu hỏi là B}
}
\end{ex}

\begin{ex} %[5H5V28]
Một mặt đồng hồ hình tròn có bán kính $15\text{ cm}$. Người ta gắn thêm một vòng viền nhôm bao quanh mép ngoài đồng hồ dày $3\text{ cm}$ (bán kính vòng ngoài là $18\text{ cm}$). Chu vi của đường viền ngoài là
\choice
{$94{,}2\text{ cm}$}
{\True $113{,}04\text{ cm}$}
{$56{,}52\text{ cm}$}
{$108\text{ cm}$}
\loigiai{
Chu vi đường viền ngoài là:
$2 \times 3{,}14 \times 18 = 113{,}04\text{ (cm)}$.\\
\textbf{Đáp án của câu hỏi là B}
}
\end{ex}

\begin{ex} %[5H5V28]
Từ một tấm tôn hình chữ nhật kích thước $40\text{ cm} \times 60\text{ cm}$, người ta cắt ra một hình tròn có đường kính lớn nhất có thể. Diện tích phần tôn bị cắt bỏ là
\choice
{$1\,140\text{ cm}^2$}
{\True $1\,144\text{ cm}^2$}
{$1\,256\text{ cm}^2$}
{$1\,200\text{ cm}^2$}
\loigiai{
Hình tròn lớn nhất có đường kính bằng chiều rộng tấm tôn, tức đường kính là $40\text{ cm}$, bán kính $r = 20\text{ cm}$.\\
Diện tích tấm tôn ban đầu là:
$40 \times 60 = 2\,400\text{ (cm}^2\text{)}$.\\
Diện tích hình tròn cắt ra là:
$3{,}14 \times 20 \times 20 = 1\,256\text{ (cm}^2\text{)}$.\\
Diện tích phần tôn cắt bỏ là:
$2\,400 - 1\,256 = 1\,144\text{ (cm}^2\text{)}$.\\
\textbf{Đáp án của câu hỏi là B}
}
\end{ex}

\begin{ex} %[5H5V28]
Trong sân trường có một bồn cây hình tròn đường kính $2\text{ m}$. Xung quanh bồn cây người ta lát một hàng gạch rộng $0{,}5\text{ m}$. Bán kính của vòng ngoài cùng của hàng gạch là
\choice
{$2{,}5\text{ m}$}
{\True $1{,}5\text{ m}$}
{$3\text{ m}$}
{$1{,}25\text{ m}$}
\loigiai{
Bán kính của bồn cây là:
$2 : 2 = 1\text{ (m)}$.\\
Bán kính vòng ngoài cùng của hàng gạch là:
$1 + 0{,}5 = 1{,}5\text{ (m)}$.\\
\textbf{Đáp án của câu hỏi là B}
}
\end{ex}

\begin{ex} %[5H5V28]
Bác thợ mộc ghép 4 tấm gỗ hình tam giác vuông bằng nhau, mỗi tấm có hai cạnh góc vuông là $30\text{ cm}$ và $40\text{ cm}$, để tạo thành một mặt bàn hình thoi. Diện tích mặt bàn hình thoi đó là
\choice
{$1\,200\text{ cm}^2$}
{\True $2\,400\text{ cm}^2$}
{$4\,800\text{ cm}^2$}
{$600\text{ cm}^2$}
\loigiai{
Diện tích của một tấm gỗ hình tam giác vuông là:
$30 \times 40 : 2 = 600\text{ (cm}^2\text{)}$.\\
Diện tích mặt bàn hình thoi là:
$600 \times 4 = 2\,400\text{ (cm}^2\text{)}$.\\
\textbf{Đáp án của câu hỏi là B}
}
\end{ex}

\begin{ex} %[5H5V28]
Một biển quảng cáo hình chữ nhật kích thước $3\text{ m} \times 6\text{ m}$. Người ta vẽ một logo hình tròn bán kính $1\text{ m}$ ở chính giữa biển. Phần diện tích còn lại của biển được sơn màu xanh. Diện tích phần sơn màu xanh là
\choice
{$15\text{ m}^2$}
{\True $14{,}86\text{ m}^2$}
{$16{,}86\text{ m}^2$}
{$17{,}14\text{ m}^2$}
\loigiai{
Diện tích biển quảng cáo là:
$3 \times 6 = 18\text{ (m}^2\text{)}$.\\
Diện tích logo hình tròn là:
$3{,}14 \times 1 \times 1 = 3{,}14\text{ (m}^2\text{)}$.\\
Diện tích phần sơn màu xanh là:
$18 - 3{,}14 = 14{,}86\text{ (m}^2\text{)}$.\\
\textbf{Đáp án của câu hỏi là B}
}
\end{ex}

\textbf{Phần 2. Trắc nghiệm đúng sai}
\setcounterpageref{ex}{0}

\begin{ex} %[5H5V28]
Bạn Lan cắt một tờ giấy hình vuông cạnh $20\text{ cm}$ thành một bộ xếp hình Tangram gồm 7 miếng ghép.
\choiceTF
{\True Diện tích ban đầu của tờ giấy hình vuông là $400\text{ cm}^2$}
{\True Tổng diện tích của cả 7 miếng ghép sau khi cắt đúng bằng $400\text{ cm}^2$}
{\True Trong bộ Tangram có hai miếng hình tam giác lớn, mỗi tam giác lớn có diện tích là $100\text{ cm}^2$}
{Diện tích của miếng hình bình hành lớn hơn diện tích của miếng hình vuông nhỏ}
\loigiai{
\begin{itemchoice}
\item Đúng. Diện tích hình vuông: $20 \times 20 = 400\text{ (cm}^2\text{)}$.
\item Đúng. Khi cắt ghép không hao hụt thì tổng diện tích giữ nguyên $400\text{ cm}^2$.
\item Đúng. Mỗi tam giác lớn chiếm $\dfrac{1}{4}$ diện tích cả bộ: $400 : 4 = 100\text{ (cm}^2\text{)}$.
\item Sai. Trong bộ Tangram chuẩn, miếng hình vuông nhỏ và miếng hình bình hành đều có diện tích bằng $\dfrac{1}{8}$ diện tích cả bộ, tức là bằng nhau ($50\text{ cm}^2$).
\end{itemchoice}
}
\end{ex}

\begin{ex} %[5H5V28]
Một sân bóng đá mini hình chữ nhật có chiều dài $25\text{ m}$ và chiều rộng $15\text{ m}$. Ở chính giữa sân có một vòng tròn tâm sân với bán kính $2\text{ m}$.
\choiceTF
{\True Diện tích toàn bộ sân bóng là $375\text{ m}^2$}
{\True Chu vi của vòng tròn giữa sân là $12{,}56\text{ m}$}
{\True Diện tích của hình tròn giữa sân là $12{,}56\text{ m}^2$}
{Diện tích phần mặt sân nằm ngoài vòng tròn giữa sân là $360\text{ m}^2$}
\loigiai{
\begin{itemchoice}
\item Đúng. Diện tích sân bóng: $25 \times 15 = 375\text{ (m}^2\text{)}$.
\item Đúng. Chu vi vòng tròn: $2 \times 3{,}14 \times 2 = 12{,}56\text{ (m)}$.
\item Đúng. Diện tích hình tròn: $3{,}14 \times 2 \times 2 = 12{,}56\text{ (m}^2\text{)}$.
\item Sai. Diện tích phần ngoài vòng tròn là: $375 - 12{,}56 = 362{,}44\text{ (m}^2\text{)} 
eq 360\text{ m}^2$.
\end{itemchoice}
}
\end{ex}

\textbf{Phần 3. Trắc nghiệm trả lời ngắn}
\setcounterpageref{ex}{0}

\begin{ex} %[5H5V28]
Một bức tranh ghép hình lá cờ hình chữ nhật có tổng diện tích bằng $600\text{ cm}^2$. Hình ngôi sao ở giữa có diện tích bằng $75\text{ cm}^2$. Phần nền đỏ còn lại của bức tranh có diện tích là bao nhiêu xăng-ti-mét vuông?
\shortans{525}
\loigiai{
Diện tích phần nền đỏ là:
$600 - 75 = 525\text{ (cm}^2\text{)}$.
}
\end{ex}

\begin{ex} %[5H5V28]
Để rào một luống hoa hình tam giác có độ dài ba cạnh lần lượt là $3{,}5\text{ m}; 4{,}2\text{ m}; 5{,}3\text{ m}$, người ta cần một đoạn lưới dài bao nhiêu đề-xi-mét?
\shortans{130}
\loigiai{
Chu vi luống hoa là:
$3{,}5 + 4{,}2 + 5{,}3 = 13\text{ (m)}$.\\
Đổi: $13\text{ m} = 130\text{ dm}$.
}
\end{ex}

\begin{ex} %[5H5V28]
Một tấm tôn hình chữ nhật kích thước $50\text{ cm} \times 80\text{ cm}$. Người ta cắt bỏ ở bốn góc bốn hình vuông bằng nhau, mỗi hình vuông có cạnh $10\text{ cm}$. Diện tích phần tôn còn lại là bao nhiêu xăng-ti-mét vuông?
\shortans{3600}
\loigiai{
Diện tích tấm tôn ban đầu là:
$50 \times 80 = 4\,000\text{ (cm}^2\text{)}$.\\
Diện tích 4 góc hình vuông cắt bỏ là:
$4 \times (10 \times 10) = 400\text{ (cm}^2\text{)}$.\\
Diện tích phần tôn còn lại là:
$4\,000 - 400 = 3\,600\text{ (cm}^2\text{)}$.
}
\end{ex}

\begin{ex} %[5H5V28]
Một cái hồ nước nhân tạo có bề mặt gồm một phần hình chữ nhật dài $30\text{ m}$, rộng $20\text{ m}$ và một phần hình tam giác tiếp giáp có cạnh đáy $20\text{ m}$, chiều cao $15\text{ m}$. Diện tích bề mặt của hồ nước đó là bao nhiêu mét vuông?
\shortans{750}
\loigiai{
Diện tích phần hình chữ nhật là:
$30 \times 20 = 600\text{ (m}^2\text{)}$.\\
Diện tích phần hình tam giác là:
$20 \times 15 : 2 = 150\text{ (m}^2\text{)}$.\\
Tổng diện tích mặt nước là:
$600 + 150 = 750\text{ (m}^2\text{)}$.
}
\end{ex}

\begin{ex} %[5H5V28]
Một bờ hồ nước có hình vuông với độ dài mỗi cạnh là $50\text{ m}$. Một người đi bộ quanh bờ hồ đúng 1 vòng thì đi được quãng đường dài bao nhiêu mét?
\shortans{200}
\loigiai{
Chu vi bờ hồ là:
$50 \times 4 = 200\text{ (m)}$.
}
\end{ex}

\begin{ex} %[5H5V28]
Người ta ghép 6 miếng gạch hình tam giác đều bằng nhau, mỗi cạnh dài $20\text{ cm}$, để tạo thành một hình lục giác đều. Chu vi của hình lục giác đều đó là bao nhiêu xăng-ti-mét?
\shortans{120}
\loigiai{
Hình lục giác đều có 6 cạnh bằng nhau, mỗi cạnh bằng độ dài cạnh của hình tam giác đều ($20\text{ cm}$).\\
Chu vi của hình lục giác đều là:
$20 \times 6 = 120\text{ (cm)}$.
}
\end{ex}
